\chapter{绪论}
\section{研究背景与意义}
近年来，无人机（UAV）集群在非结构化环境中的应用前景广阔，如灾害救援、环境监测和物资协同配送等。然而，在复杂动态环境中，传统的集中式控制方法面临维度灾难和通信延迟等瓶颈。

\section{国内外研究现状}
\subsection{多无人机协同规划与避障方法}
现有方法多依赖于精确的环境地图，在未知或高动态环境下面临鲁棒性不足的问题。

\subsection{深度强化学习（DRL）研究进展}
深度强化学习（DRL）为解决高维状态空间的决策问题提供了有效途径，但在多尺度耦合的物理系统中，直接应用端到端强化学习往往会导致灾难性遗忘和控制高频抖动。

\subsection{自抗扰控制（LADRC）及其应用研究现状}
线性自抗扰控制（LADRC）具有不依赖精确模型的优势，但其固定参数在面对剧烈风场扰动时难以保持最优性能。

\section{本文研究内容与创新点}
本文提出一种分层引导课程强化学习（HGC-RL）架构，主要创新点包括：
\begin{itemize}
    \item 提出解耦制导与控制的分层异步优化框架。
    \item 设计时序增强（TSA）的自适应 LADRC 底层控制机制。
    \item 构建基于金字塔优先经验池（PHM）和课程学习的多机协同策略。
\end{itemize}

\section{论文结构安排}
本文共分为五章，第一章介绍研究背景；第二章建立问题模型；第三章设计底层运动控制；第四章设计上层规划策略并进行仿真验证；第五章总结全文。
